\par\clearpage
\SourceLeaf{132}\GaussVTheoryEarlyFolio{128}
\[
V=-\int\frac{\mathrm{d}\mu}{\rho}
\]
Es hat also $V$ in jedem Punkte des Raumes einen bestimmten Werth, oder es ist eine Function von $x,y,z$, oder auch von je drei andern veränderlichen Grössen, wodurch man die Punkte des Raumes unterscheidet. Die magnetische Kraft $\psi$ in jedem Punkte des Raumes, und die Componenten $\xi,\eta,\zeta$, die aus der Zerlegung von $\psi$ parallel mit den Coordinatenaxen entstehen, finden sich dann durch die Formeln
\[
\xi=\frac{\mathrm{d}V}{\mathrm{d}x},\quad \eta=\frac{\mathrm{d}V}{\mathrm{d}y},\quad \zeta=\frac{\mathrm{d}V}{\mathrm{d}z},\quad \psi=\BareRoot{(\xi\xi+\eta\eta+\zeta\zeta)}
\]

\Article{5}
Es sollen nun zuvörderst einige allgemeine von der Form der Function $V$ unabhängige Sätze entwickelt werden, die wegen ihrer Einfachheit und Eleganz merkwürdig sind.

Das vollständige Differential von $V$ wird
\[
\begin{aligned}
\mathrm{d}V&=\frac{\mathrm{d}V}{\mathrm{d}x}\cdot\mathrm{d}x+\frac{\mathrm{d}V}{\mathrm{d}y}\cdot\mathrm{d}y+\frac{\mathrm{d}V}{\mathrm{d}z}\cdot\mathrm{d}z\\
&=\xi\,\mathrm{d}x+\eta\,\mathrm{d}y+\zeta\,\mathrm{d}z
\end{aligned}
\]

Bezeichnet man mit $\mathrm{d}s$ die Entfernung zwischen den beiden Punkten, auf welche sich $V$ und $V+\mathrm{d}V$ beziehen, und mit $\theta$ den Winkel, welchen die Richtung der magnetischen Kraft $\psi$ mit $\mathrm{d}s$ macht, so wird
\[
\mathrm{d}V=\psi\cos\theta\cdot\mathrm{d}s
\]
weil $\frac{\xi}{\psi},\frac{\eta}{\psi},\frac{\zeta}{\psi}$ die Cosinus der Winkel sind, welche die Richtung von $\psi$ mit den Coordinatenaxen macht, hingegen $\frac{\mathrm{d}x}{\mathrm{d}s},\frac{\mathrm{d}y}{\mathrm{d}s},\frac{\mathrm{d}z}{\mathrm{d}s}$ die Cosinus der Winkel zwischen $\mathrm{d}s$ und denselben Axen. Es ist also $\frac{\mathrm{d}V}{\mathrm{d}s}$ gleich der auf die Richtung von $\mathrm{d}s$ projicirten Kraft; dasselbe folgt auch schon aus der Gleichung $\frac{\mathrm{d}V}{\mathrm{d}x}=\xi$, wenn man sich erinnert, dass die Coordinatenaxen nach Willkür gewählt werden können.

\Article{6}
Werden zwei Punkte im Raume, $P^0,P'$ durch eine beliebige Linie verbunden, wovon $\mathrm{d}s$ ein unbestimmtes Element vorstellt, und bedeutet, wie vorhin $\theta$ den Winkel zwischen $\mathrm{d}s$ und der Richtung der daselbst Statt findenden
\par\clearpage
\SourceLeaf{133}\GaussVTheoryEarlyFolio{129}
\noindent magnetischen Kraft, und $\psi$ deren Intensität, so ist
\[
\int\psi\cos\theta\cdot\mathrm{d}s=V'-V^0
\]
wenn man die Integration durch die ganze Linie ausdehnt, und mit $V^0,V'$ die Werthe von $V$ an den Endpunkten bezeichnet.

Folgende Corollarien dieses fruchtbaren Satzes verdienen hier besonders angeführt zu werden:

I. Das Integral $\int\psi\cos\theta\cdot\mathrm{d}s$ behält einerlei Werth: auf welchem Wege man auch von $P^0$ nach $P'$ übergeht.

II. Das Integral $\int\psi\cos\theta\cdot\mathrm{d}s$ durch die ganze Länge irgend einer in sich zurückkehrenden Linie ausgedehnt, ist immer $=0$.

III. In einer geschlossenen Linie muss, wenn nicht durchgehends $\theta=90^0$ ist, ein Theil der Werthe von $\theta$ kleiner und ein Theil grösser als $90^0$ sein.

\Article{7}
\expandafter\def\csname HistoricalNoteText15\endcsname{Konnte die Function $V$ jede willkürlich aufgestellte Form haben, so konnte in besondern Fällen ein Maximum- oder Minimum-Werth von $V$ einem isolirten Punkte oder einer isolirten Linie entsprechen, um welchen oder um welche ringsum bloss kleinere oder bloss grössere Werthe Statt finden würden, oder auch einer Fläche, auf deren \textit{beiden} Seiten zugleich kleinere oder grössere Werthe gälten. Allein die Bedingungen, denen die Function $V$ unterworfen ist, lassen diese Ausnahmsfälle nicht zu. Eine ausführliche Entwickelung dieses Gegenstandes muss aber, da sie für unsern gegenwärtigen Zweck unnöthig ist, einer andern Gelegenheit vorbehalten bleiben.}
Die Fläche, in deren sämmtlichen Punkten $V$ einerlei bestimmten Werth $=V^0$ hat, scheidet die Punkte des Raumes, in welchen $V$ einen Werth grösser als $V^0$ hat, von denen, wo der Werth kleiner als $V^0$ ist\HistoricalNote{15}. Aus dem Satz des Art. 5. folgt leicht, dass die magnetische Kraft in jedem Punkte dieser Fläche eine gegen die Fläche senkrechte Richtung hat, und zwar nach der Seite zu, auf welcher die grössern Werthe von $V$ Statt finden. Ist $\mathrm{d}s$ eine unendlich kleine gegen die Fläche senkrechte Linie, und $V^0+\mathrm{d}V^0$ der Werth von $V$ an dem andern Endpunkte derselben, so wird die Intensität der magnetischen Kraft $=\frac{\mathrm{d}V^0}{\mathrm{d}s}$. Die Gesammtheit der Punkte, wofür $V=V^0+\mathrm{d}V^0$ ist, bildet eine zweite der ersten unendlich nahe Fläche, und an den verschiedenen Stellen des ganzen Zwischenraumes ist die Intensität der magnetischen Kraft der Entfernung beider Flächen von einander verkehrt proportional. Lässt man $V$ durch unendlich kleine aber gleiche Stufen sich ändern, so entsteht dadurch ein System
\par\GaussVTheoryGathering{17}
\par\clearpage
\SourceLeaf{134}\GaussVTheoryEarlyFolio{130}
\noindent von Flächen, die den Raum in unendlich dünne Schichten abtheilen, und die verkehrte Proportionalität der Dicke der Schichten zu der Intensität der magnetischen Kraft gilt dann nicht bloss für verschiedene Stellen einer und derselben Schicht, sondern auch für verschiedene Schichten.

\Article{8}
Wir wollen nun das Verhalten der Werthe von $V$ auf der Oberfläche der Erde betrachten.

Es sei in einem Punkte $P$ der Erdoberfläche $\psi$ die Intensität, $PM$ die Richtung der ganzen magnetischen Kraft; $\omega$ die Intensität, $PN$ die Richtung der auf die horizontale Ebene projicirten Kraft, oder $PN$ die Richtung des magnetischen Meridians, in dem Sinn vom Südpol der Magnetnadel zum Nordpol; $i$ der Winkel zwischen $PM$ und $PN$ oder die Inclination; $\theta,t$ die Winkel zwischen dem Elemente $\mathrm{d}s$ einer auf der Erdoberfläche liegenden Linie und den Richtungen $PM,PN$; endlich entsprechen $V$ und $V+\mathrm{d}V$ dem Anfangs- und Endpunkte von $\mathrm{d}s$. Wir haben folglich
\[
\cos\theta=\cos i\cos t,\qquad \omega=\psi\cos i
\]
und die Gleichung des Art. 5. verwandelt sich in
\[
\mathrm{d}V=\omega\cos t\cdot\mathrm{d}s
\]

Sind also zwei Punkte $P^0,P'$ auf der Erdoberfläche, in welchen $V$ die Werthe $V^0,V'$ hat, durch eine ganz auf der Erdoberfläche liegende Linie verbunden, von welcher $\mathrm{d}s$ ein unbestimmtes Element bedeutet, so ist
\[
\int\omega\cos t\cdot\mathrm{d}s=V'-V^0
\]
wenn die Integration durch die ganze Linie ausgedehnt wird, und offenbar gelten nun auch hier drei den im Art. 6. angeführten ganz ähnliche Corollarien, nemlich:

I. Das Integral $\int\omega\cos t\,\mathrm{d}s$ behält einerlei Werth, auf welchem Wege auf der Oberfläche der Erde man auch von $P^0$ nach $P'$ übergeht.

II. Das Integral $\int\omega\cos t\cdot\mathrm{d}s$ durch die ganze Länge einer auf der Oberfläche der Erde liegenden geschlossenen Linie ist immer $=0$.

III. In einer solchen geschlossenen Linie muss nothwendig, falls nicht
\par\clearpage
\SourceLeaf{135}\GaussVTheoryEarlyFolio{131}
\noindent durchgehends $t=90^0$ ist. ein Theil der Werthe von $t$ spitz und ein Theil stumpf sein.

\Article{9}
Die Sätze I. und II. des vorhergehenden Artikels (welche eigentlich nur zwei verschiedene Einkleidungen derselben Sache sind) lassen sich wenigstens näherungsweise an wirklichen Beobachtungen prüfen. Es sei $P^0P'P''....P^0$ ein Polygon auf der Erdoberfläche, dessen Seiten die kürzesten Linien zwischen ihren Endpunkten, also, wenn man die Erde hier nur als kugelförmig betrachtet, grösste Kreisbögen sind. Es seien $\omega^0,\omega',\omega''$ u. s. w. die Intensitäten der horizontalen magnetischen Kraft in den Punkten $P^0,P',P''$ u. s. w.; ferner $\delta^0,\delta',\delta''$ u. s. w. die Declinationen, die man nach üblicher Weise westlich vom Nordpunkte als positiv, östlich als negativ betrachten mag: endlich sei $(01)$ das Azimuth der Linie $P^0P'$ in $P^0$, und zwar nach üblicher Weise von Süden aus nach Westen herumgezählt; eben so $(10)$ das Azimuth derselben Linie rückwärts genommen in $P'$ u. s. w.

Man bemerke. dass $t$ zwar in jeder Polygonseite sich nach der Stetigkeit ändert, in den Eckpunkten hingegen sprungsweise, und also in diesen zwei verschiedene Werthe hat; z. B. in $P'$ hat $t$

den Werth $(10)+\delta'$, insofern $P'$ der Endpunkt von $P^0P'$ ist,

den Werth $180^0+(12)+\delta'$, insofern $P'$ der Anfangspunkt von $P'P''$ ist.

Von dem Integral $\int\omega\cos t\cdot\mathrm{d}s$, durch $P^0P'$ ausgedehnt, kann man als genäherten Werth betrachten
\[
\tfrac12(\omega^0\cos t^0+\omega'\cos t')\cdot P^0P'
\]
wenn $t^0,t'$ die Werthe von $t$ in $P^0$ als Anfangspunkt und in $P'$ als Endpunkt von $P^0P'$ bedeuten; diese Annäherung ist alles, was man erlangen kann, insofern man die Werthe von $\omega$ und $t$ eben nur in den Endpunkten $P^0,P'$ hat, und sie ist um so zulässiger, je kleiner die Linie ist. Der angegebene Ausdruck ist in unsern Bezeichnungen
\[
=\tfrac12\{\omega'\cos((10)+\delta')-\omega^0\cos((01)+\delta^0)\}\cdot P^0P'
\]
Auf ähnliche Art ist der genäherte Werth des Integrals, durch $P'P''$ ausgedehnt,
\par\GaussVTheoryGathering{17\textsuperscript{*}}
\par\clearpage
\SourceLeaf{136}\GaussVTheoryEarlyFolio{132}
\[
=\tfrac12\{\omega''\cos((21)+\delta'')-\omega'\cos((12)+\delta')\}\cdot P'P''
\]
u. s. f. durch das ganze Polygon.

Für ein Dreieck gibt also unser Satz die näherungsweise richtige Gleichung
\[
\begin{aligned}
&\omega^0\{P^0P'\cos((01)+\delta^0)-P^0P''\cos((02)+\delta^0)\}\\
{}+{}&\omega'\{P'P''\cos((12)+\delta')-P^0P'\cos((10)+\delta')\}\\
{}+{}&\omega''\{P^0P''\cos((20)+\delta'')-P'P''\cos((21)+\delta'')\}\\
={}&0
\end{aligned}
\]

Offenbar sind bei dieser Gleichung die Einheiten für die Intensitäten und Distanzen willkürlich.

\Article{10}
Als ein Beispiel wollen wir die Formel auf die magnetischen Elemente von
\begin{center}
\begin{tabular}{l|l|l|l}
Göttingen&$\delta^0=18^0\ 38'$&$i^0=67^0\ 56'$&$\psi^0=1{,}357$\\
Mailand&$\delta'=18\ 33$&$i'=63\ 49$&$\psi'=1{,}294$\\
Paris&$\delta''=22\ 4$&$i''=67\ 24$&$\psi''=1{,}348$
\end{tabular}
\end{center}
anwenden, woraus
\[
\begin{aligned}
\omega^0&=0{,}50980\\
\omega'&=0{,}57094\\
\omega''&=0{,}51804
\end{aligned}
\]
folgt. Legt man die geographische Lage
\begin{center}
\begin{tabular}{l|l|l}
Göttingen&$51^0\ 32'$ Breite&$9^0\ 58'$ Länge von Greenwich\\
Mailand&$45\ 28$&$9\ 9$\\
Paris&$48\ 52$&$2\ 21$
\end{tabular}
\end{center}
zum Grunde, und führt die Rechnung nur wie auf der Kugelfläche, so findet sich
\[
\begin{aligned}
\left.\begin{aligned}(01)&=5^0\ 11'\ 31''\\(10)&=184\ 35\ 35\end{aligned}\right\}&\quad P^0P'=6^0\ 5'\ 20''\\[4pt]
\left.\begin{aligned}(12)&=128\ 47\ 31\\(21)&=303\ 48\ 1\end{aligned}\right\}&\quad P'P''=5\ 44\ 6\\[4pt]
\left.\begin{aligned}(20)&=238\ 20\ 20\\(02)&=64\ 10\ 12\end{aligned}\right\}&\quad P^0P''=5\ 32\ 4
\end{aligned}
\]
